% !TEX program = xelatex
% =====================================================================
%  SLIDE BÁO CÁO ĐỒ ÁN TRÍ TUỆ NHÂN TẠO, HỌC PHẦN AIP202
%  Build:  latexmk -xelatex slide.tex   (Overleaf: Compiler = XeLaTeX)
%  Ảnh lấy từ ../bao-cao-typhoon-vn/figs; thiếu ảnh thì slide vẫn build.
% =====================================================================
\documentclass[aspectratio=169, 11pt]{beamer}

\usepackage{fontspec}
\usepackage{booktabs}
\usepackage{tikz}
\usetikzlibrary{arrows.meta, positioning}

% ---------------------------------------------------------------------
%  Thông tin nhóm: sửa ở đây
% ---------------------------------------------------------------------
\newcommand{\sinhvien}{Nguyễn Viết Mạnh (2351220241) \and Nguyễn Đình Phú (2351220247)}
\newcommand{\gvhd}{TS. Đỗ Phúc Hảo}
\newcommand{\lop}{23CT5}

% ---------------------------------------------------------------------
%  Nhận diện DAU: cùng bảng màu với báo cáo
% ---------------------------------------------------------------------
\definecolor{daumaroon}{HTML}{990000}
\definecolor{daugold}{HTML}{FBAE40}
\definecolor{daunavy}{HTML}{0E2841}
\definecolor{gray100}{HTML}{F3F4F6}

\usetheme{Madrid}
\usecolortheme[named=daunavy]{structure}
\setbeamercolor{title}{fg=white, bg=daunavy}
\setbeamercolor{frametitle}{fg=white, bg=daunavy}
\setbeamercolor{author in head/foot}{fg=white, bg=daumaroon}
\setbeamercolor{title in head/foot}{fg=white, bg=daunavy}
\setbeamercolor{date in head/foot}{fg=daunavy, bg=daugold}
\setbeamercolor{block title}{fg=white, bg=daumaroon}
\setbeamercolor{block body}{bg=gray100}
\setbeamercolor{alerted text}{fg=daumaroon}
\setbeamertemplate{navigation symbols}{}
\setbeamertemplate{itemize items}[square]
\setbeamerfont{frametitle}{series=\bfseries}

% Chèn ảnh nếu có, nếu không thì in khung báo thiếu để slide vẫn build được
\graphicspath{{../bao-cao-typhoon-vn/figs/}}
\newcommand{\hinh}[2][]{%
  \IfFileExists{../bao-cao-typhoon-vn/figs/#2}%
    {\includegraphics[#1]{#2}}%
    {\fbox{\parbox{0.6\linewidth}{\centering\small Thiếu ảnh \texttt{\detokenize{#2}}}}}}

\tikzset{
  khoi/.style = {rectangle, rounded corners=2pt, draw=daumaroon, line width=0.8pt,
                 fill=white, align=center, font=\scriptsize, minimum height=9mm,
                 inner sep=3pt},
  kho/.style  = {khoi, draw=daunavy, fill=gray100},
  mui/.style  = {-{Stealth[length=2mm]}, line width=0.8pt, color=daumaroon},
}

\title[Gợi ý đường đi của bão]{Xây dựng hệ thống hỗ trợ gợi ý\\ đường đi của cơn bão Việt Nam}
\subtitle{Báo cáo đồ án Trí tuệ nhân tạo \textperiodcentered{} Học phần AIP202}
\author[Nhóm \lop]{\sinhvien}
\institute[DAU]{Khoa Công nghệ Thông tin, Trường Đại học Kiến trúc Đà Nẵng\\[2pt]
  Giảng viên hướng dẫn: \gvhd}
\date{Năm học 2025--2026}

\AtBeginSection[]{%
  \begin{frame}{Nội dung}
    \tableofcontents[currentsection]
  \end{frame}}

\begin{document}

\begin{frame}[plain]
  \titlepage
\end{frame}

% =====================================================================
\section{Bài toán và mục tiêu}
% =====================================================================

\begin{frame}{Vì sao chọn đề tài này}
  \begin{itemize}
    \item Việt Nam nằm cạnh Biển Đông, vùng xoáy thuận nhiệt đới hoạt động mạnh.
    \item Câu hỏi quan trọng nhất khi phòng chống bão:
          \alert{tâm bão sẽ ở đâu trong 1 đến 3 ngày tới?}
    \item Dự báo nghiệp vụ cần mô hình số trị và siêu máy tính.
    \item Hơn 40 năm dữ liệu quỹ đạo bão đã được công bố công khai.
  \end{itemize}
  \vspace{0.6em}
  \begin{block}{Hai câu hỏi của đồ án}
    \begin{enumerate}
      \item AI chỉ học từ dữ liệu lịch sử gợi ý đường đi chính xác đến mức nào?
      \item Lập luận theo tình huống, học máy và học sâu khác nhau ra sao trên
            cùng một bài toán?
    \end{enumerate}
  \end{block}
\end{frame}

\begin{frame}{Phát biểu bài toán}
  \begin{columns}[T]
    \begin{column}{0.48\textwidth}
      \begin{block}{Đầu vào}
        4 quan trắc gần nhất, cách nhau 6 giờ (24 giờ lịch sử): thời điểm,
        vĩ độ, kinh độ, gió, khí áp, cấp cường độ.
      \end{block}
    \end{column}
    \begin{column}{0.48\textwidth}
      \begin{block}{Đầu ra}
        Vị trí tâm bão tại +6, +12, +24, +48, +72 giờ; nón bất định 80\%;
        5 cơn bão lịch sử tương tự nhất.
      \end{block}
    \end{column}
  \end{columns}
  \vspace{0.8em}
  \textbf{Phạm vi:} Tây Bắc Thái Bình Dương, 0--30°N, 100--140°E, 1980--2024.

  \textbf{Không làm:} dự báo cường độ, mưa, lũ, nước dâng; không phát cảnh báo
  chính thức, không thay thế bản tin của NCHMF.
\end{frame}

\begin{frame}{Năm mục tiêu có thể kiểm chứng}
  \begin{description}[MT5]
    \item[MT1] Quy trình dữ liệu từ NOAA IBTrACS, chia tập theo năm, không rò rỉ.
    \item[MT2] So sánh ít nhất 4 phương pháp AI với 2 mốc, có kiểm định thống kê.
    \item[MT3] Giảm sai số 72 giờ ít nhất 20\% so với persistence.
    \item[MT4] Nón bất định có độ phủ thực tế xấp xỉ 80\%.
    \item[MT5] Triển khai thành API và bản đồ, có giải thích bằng bão tương tự.
  \end{description}
\end{frame}

% =====================================================================
\section{Thiết kế giải pháp}
% =====================================================================

\begin{frame}{Kiến trúc tổng thể}
  \centering
  \resizebox{\ifdim\width>\textwidth\textwidth\else\width\fi}{!}{%
  \begin{tikzpicture}[node distance=7mm and 5mm]
    \node[kho]  (data)  {IBTrACS WP\\1980--2024};
    \node[khoi, right=of data]  (clean) {Làm sạch,\\chuẩn hóa};
    \node[khoi, right=of clean] (feat)  {Tạo đặc trưng\\(47 + 11 biến)};
    \node[khoi, right=of feat]  (split) {Cửa sổ chuỗi,\\chia theo năm};
    \node[khoi, right=of split] (train) {Huấn luyện,\\chọn trên val};
    \node[kho,  right=of train] (bundle){Gói mô hình\\+ analog};
    \node[kho,  below=of feat]  (env)   {NCEP/NCAR R1\\(dòng dẫn đường)};
    \node[khoi, below=of bundle](api)   {FastAPI\\dự báo + nón};
    \node[kho,  left=of api]    (web)   {Bản đồ web\\+ bão tương tự};
    \draw[mui] (data) -- (clean);
    \draw[mui] (clean) -- (feat);
    \draw[mui] (feat) -- (split);
    \draw[mui] (split) -- (train);
    \draw[mui] (train) -- (bundle);
    \draw[mui] (env) -- (feat);
    \draw[mui] (bundle) -- (api);
    \draw[mui] (api) -- (web);
  \end{tikzpicture}}

  \vspace{0.8em}
  \raggedright
  Huấn luyện chạy ngoại tuyến, phục vụ chạy trực tuyến. Hai nửa chỉ giao tiếp
  qua \alert{gói mô hình có mã băm SHA-256}, nên huấn luyện lại không phải sửa
  phần phục vụ.
\end{frame}

\begin{frame}{Dữ liệu và chia tập không rò rỉ}
  Sau làm sạch: \alert{24.024 quan trắc của 1.089 cơn bão}.
  \vspace{0.4em}

  \centering
  \begin{tabular}{llrl}
    \toprule
    \textbf{Tập} & \textbf{Giai đoạn} & \textbf{Số bão} & \textbf{Mục đích} \\
    \midrule
    Huấn luyện   & 1980--2016 & 895 & Học tham số (khoảng 7.800 cửa sổ) \\
    Validation   & 2017--2020 & 110 & Chọn mô hình, siêu tham số, bán kính nón \\
    Kiểm thử     & 2021--2024 & 84  & Đánh giá một lần (714 cửa sổ) \\
    Ảnh hưởng VN & 2021--2024 & 27  & Tập con của kiểm thử (278 cửa sổ) \\
    \bottomrule
  \end{tabular}

  \vspace{0.6em}
  \raggedright
  \begin{itemize}
    \item Chia \textbf{theo năm, giữ nguyên cả cơn bão}: quan trắc của một cơn
          bão không nằm ở hai tập.
    \item Bộ chuẩn hóa chỉ fit trên tập huấn luyện.
  \end{itemize}
\end{frame}

\begin{frame}{Đặc trưng và cách biểu diễn đầu ra}
  \begin{columns}[T]
    \begin{column}{0.52\textwidth}
      \textbf{47 đặc trưng mỗi quan trắc}
      \begin{itemize}
        \item Vị trí và độ trễ 1--3 bước
        \item Động học: tốc độ, hướng, góc rẽ
        \item Thời gian: sin/cos giờ, tháng, mùa bão
        \item Địa lý: khoảng cách tới bờ biển Việt Nam
        \item Cường độ: gió, khí áp, cấp bão
      \end{itemize}
      \textbf{Mở rộng:} 11 đặc trưng dòng dẫn đường khí quyển.
    \end{column}
    \begin{column}{0.44\textwidth}
      \begin{block}{Học độ dịch chuyển}
        \[
          \hat{\mathbf{p}}_{t+h} = \mathbf{p}_t + \Delta\hat{\mathbf{p}}_h
        \]
        Vị trí hiện tại đã biết chính xác. Mạng chỉ cần học hướng và tốc độ
        di chuyển.
      \end{block}
    \end{column}
  \end{columns}
\end{frame}

\begin{frame}{Bảy phương pháp thuộc bốn nhóm}
  \centering
  \begin{tabular}{lll}
    \toprule
    \textbf{Nhóm} & \textbf{Phương pháp} & \textbf{Cấu hình} \\
    \midrule
    Mốc so sánh & Persistence       & Giữ nguyên vận tốc 6 giờ gần nhất \\
    Mốc so sánh & CLIPER-proxy      & Persistence có suy giảm \\
    AI cổ điển  & k-NN analog       & $k = 40$, có giải thích bằng bão tương tự \\
    Học máy     & Gradient boosting & Hồi quy độ dịch chuyển \\
    Học sâu     & LSTM              & 2 tầng, 128 nơ-ron ẩn, dropout 0,2 \\
    Học sâu     & Attention-LSTM    & LSTM + attention cộng \\
    Học sâu     & Transformer       & Encoder self-attention \\
    \bottomrule
  \end{tabular}

  \vspace{0.6em}
  \raggedright
  Mọi lựa chọn chỉ dựa trên validation 2017--2020. Tập kiểm thử được chạy
  \alert{một lần} sau khi chốt cấu hình. Mỗi mạng chạy 3 seed.
\end{frame}

% =====================================================================
\section{Kết quả và đánh giá}
% =====================================================================

\begin{frame}{Sai số trung bình trên tập kiểm thử 2021--2024 (km)}
  \centering
  \begin{tabular}{lrrrrr}
    \toprule
    \textbf{Phương pháp} & \textbf{6h} & \textbf{12h} & \textbf{24h} & \textbf{48h} & \textbf{72h} \\
    \midrule
    Persistence       & 31 & 68 & 155 & 390 & 676 \\
    CLIPER-proxy      & 30 & 65 & 151 & 388 & 677 \\
    k-NN analog       & 32 & 64 & 133 & 290 & \textbf{462} \\
    Gradient boosting & \textbf{29} & \textbf{60} & \textbf{129} & 292 & 476 \\
    LSTM              & 33 & 65 & 133 & \textbf{289} & 464 \\
    Attention-LSTM    & 35 & 68 & 138 & 297 & 473 \\
    Transformer       & 33 & 64 & 133 & 293 & 468 \\
    \bottomrule
  \end{tabular}

  \vspace{0.6em}
  \raggedright
  84 cơn bão, 714 cửa sổ. Tại 72 giờ: \alert{676~km xuống 462~km, giảm 32\%}
  (mục tiêu MT3 là 20\%).
\end{frame}

\begin{frame}{Sai số theo thời hạn dự báo}
  \begin{columns}[c]
    \begin{column}{0.58\textwidth}
      \centering
      \hinh[width=\linewidth]{error_by_horizon.png}
    \end{column}
    \begin{column}{0.38\textwidth}
      \begin{itemize}
        \item Ở 6 giờ các đường gần như trùng nhau: persistence gần tối ưu.
        \item Từ 24 giờ, nhóm phương pháp học tách khỏi hai mốc.
        \item Không phương pháp học nào dẫn đầu ở mọi mốc.
      \end{itemize}
    \end{column}
  \end{columns}
\end{frame}

\begin{frame}{Kiểm định thống kê và nón bất định}
  \begin{columns}[T]
    \begin{column}{0.54\textwidth}
      \textbf{Chênh lệch 72 giờ so với persistence (km)}\\[2pt]
      \small
      \begin{tabular}{lr}
        \toprule
        \textbf{Phương pháp} & \textbf{Khoảng tin cậy 95\%} \\
        \midrule
        k-NN analog       & $-201$ \ [$-268$; $-131$] \\
        Gradient boosting & $-191$ \ [$-252$; $-126$] \\
        LSTM              & $-203$ \ [$-262$; $-145$] \\
        Attention-LSTM    & $-173$ \ [$-229$; $-116$] \\
        Transformer       & $-200$ \ [$-257$; $-142$] \\
        \bottomrule
      \end{tabular}

      \vspace{0.3em}
      Bootstrap theo cơn bão, 50 cơn bão, 5.000 lần lấy mẫu lại.
    \end{column}
    \begin{column}{0.42\textwidth}
      \begin{block}{Nón bất định 80\%}
        Bán kính: 52, 97, 210, 455, 714~km tại 6, 12, 24, 48, 72 giờ.

        Độ phủ thực tế trên tập kiểm thử: \alert{82--86\%}.
      \end{block}
      Mọi phương pháp học đều tốt hơn persistence có ý nghĩa thống kê; giữa
      LSTM, Transformer, k-NN và gradient boosting không có khác biệt có ý
      nghĩa.
    \end{column}
  \end{columns}
\end{frame}

\begin{frame}{Hai thí nghiệm bổ sung}
  \begin{columns}[T]
    \begin{column}{0.46\textwidth}
      \textbf{Học độ dịch chuyển (LSTM, km)}\\[2pt]
      \small
      \begin{tabular}{lrrr}
        \toprule
        \textbf{Mục tiêu học} & \textbf{6h} & \textbf{24h} & \textbf{72h} \\
        \midrule
        Tọa độ tuyệt đối & 115 & 164 & 473 \\
        Độ dịch chuyển   & \textbf{33} & \textbf{133} & \textbf{464} \\
        \bottomrule
      \end{tabular}

      \vspace{0.4em}
      \normalsize
      Sai số 6 giờ giảm khoảng 3,5 lần.
    \end{column}
    \begin{column}{0.50\textwidth}
      \textbf{Thêm dòng dẫn đường (km)}\\[2pt]
      \small
      \begin{tabular}{lrrr}
        \toprule
        \textbf{Cấu hình} & \textbf{24h} & \textbf{72h} & \textbf{TB} \\
        \midrule
        LSTM chỉ quỹ đạo        & 131 & 469 & 197 \\
        LSTM + dẫn đường        & 118 & 465 & 188 \\
        Transformer + dẫn đường & 115 & 450 & 183 \\
        Kết hợp hai mô hình     & \textbf{112} & \textbf{443} & \textbf{179} \\
        \bottomrule
      \end{tabular}

      \vspace{0.4em}
      \normalsize
      Giảm 35~km tại 72 giờ, khoảng tin cậy [$-62$; $-9$].
    \end{column}
  \end{columns}
  \vspace{0.6em}
  \begin{block}{Kết luận chính}
    Đổi kiến trúc hay siêu tham số không giúp gì; thêm thông tin khí quyển thì
    có. \alert{Nút thắt là thông tin đầu vào, không phải kiến trúc mô hình.}
  \end{block}
\end{frame}

\begin{frame}{Nghiên cứu trường hợp: bão YAGI 2024}
  \begin{columns}[c]
    \begin{column}{0.56\textwidth}
      \centering
      \hinh[width=\linewidth, height=0.7\textheight, keepaspectratio]{case_WP122024.png}
    \end{column}
    \begin{column}{0.40\textwidth}
      Phát dự báo lúc 06Z 04/09/2024, sai số 72 giờ:
      \begin{itemize}
        \item Gradient boosting: \textbf{407~km}
        \item k-NN analog: 419~km
        \item LSTM: 541~km
        \item Transformer: 678~km
        \item Persistence: 701~km
      \end{itemize}
      Phần lớn mô hình gợi ý bão chuyển hướng lên phía bắc, trong khi YAGI đi
      thẳng về phía tây.
    \end{column}
  \end{columns}
\end{frame}

% =====================================================================
\section{Hệ thống và kết luận}
% =====================================================================

\begin{frame}{Hệ thống đã triển khai}
  \begin{columns}[T]
    \begin{column}{0.48\textwidth}
      \textbf{Phần phục vụ}
      \begin{itemize}
        \item API FastAPI: dự báo, nón, bão đang hoạt động
        \item Gói mô hình kiểm tra mã băm trước khi nạp
        \item Phản hồi kèm phiên bản mô hình, nguồn dữ liệu, cảnh báo
        \item Xác thực bằng khóa API, giới hạn tần suất
      \end{itemize}
    \end{column}
    \begin{column}{0.48\textwidth}
      \textbf{Giao diện và kiểm thử}
      \begin{itemize}
        \item Bản đồ web: quỹ đạo, nón, bão tương tự
        \item Demo công khai: \texttt{baoviet2026.pages.dev}
        \item \alert{177 ca kiểm thử Python đạt}, 4 ca bỏ qua
        \item 42 ca kiểm thử bản đồ web đạt
      \end{itemize}
    \end{column}
  \end{columns}
  \vspace{0.8em}
  \begin{block}{Lưu ý}
    Hệ thống chỉ phục vụ học tập và nghiên cứu, không phải hệ thống cảnh báo
    chính thức.
  \end{block}
\end{frame}

\begin{frame}{Đối chiếu mục tiêu}
  \centering
  \begin{tabular}{lll}
    \toprule
    \textbf{Mục tiêu} & \textbf{Kết quả} & \textbf{Mức độ} \\
    \midrule
    MT1 & 24.024 quan trắc, chia theo năm, scaler fit trên train & Đạt \\
    MT2 & 7 phương pháp, 3 seed, bootstrap theo cơn bão & Đạt \\
    MT3 & 72 giờ: 676~km xuống 462~km, giảm 32\% & Đạt \\
    MT4 & Nón 80\% phủ 82--86\% vị trí thực & Đạt \\
    MT5 & API, bản đồ web, demo công khai & Đạt một phần \\
    \bottomrule
  \end{tabular}

  \vspace{0.8em}
  \raggedright
  MT5 đạt một phần vì mô hình đang phục vụ là LSTM chỉ dùng quỹ đạo, chưa dùng
  dòng dẫn đường.
\end{frame}

\begin{frame}{Hạn chế và hướng phát triển}
  \begin{columns}[T]
    \begin{column}{0.48\textwidth}
      \textbf{Hạn chế}
      \begin{itemize}
        \item Mô hình có dòng dẫn đường chưa vận hành được vì cần dữ liệu khí
              quyển tại thời điểm dự báo.
        \item YAGI 2024 vẫn sai 620--670~km tại 72 giờ dù có dòng dẫn đường.
        \item Nón là đường tròn đẳng hướng.
        \item Chỉ một cách chia theo năm, chưa kiểm định chéo theo mùa.
        \item Chưa nghiệm thu trong môi trường vận hành thật.
      \end{itemize}
    \end{column}
    \begin{column}{0.48\textwidth}
      \textbf{Hướng phát triển}
      \begin{itemize}
        \item Nối dữ liệu phân tích khí quyển thời gian thực (ví dụ GFS).
        \item Nón bất định dạng elip theo hướng di chuyển.
        \item Kiểm định chéo theo mùa bão.
        \item Nguồn bão đang hoạt động ít trễ hơn.
      \end{itemize}
    \end{column}
  \end{columns}
\end{frame}

\begin{frame}[plain]
  \centering
  \vfill
  {\Huge\bfseries\color{daunavy} Xin cảm ơn quý thầy cô và các bạn}

  \vspace{1.2em}
  {\large Nhóm xin lắng nghe câu hỏi và góp ý.}
  \vfill
\end{frame}

\end{document}
